**Note on the Research Program: The Latent Cognition Trilogy**

This article constitutes the first component of a structured theoretical research program aimed at formalizing internal mechanisms of machine cognition. The program, titled *The Latent Cognition Trilogy*, is organized in three stages: (i) the definition of latent memory as a differentiable field within neural architectures, (ii) the emergence of symbolic structure under compression and curriculum constraints, and (iii) the architectural organization of computation through hierarchical routing mechanisms.

Rather than presenting a unified theory, each volume isolates a specific class of mechanisms and develops corresponding formal descriptions and testable propositions. The present work (Volume I) introduces Differentiable Latent Indexing (DLI) as a candidate formalism for internal memory retrieval within transformer-based models.

Subsequent volumes extend this foundation by examining the conditions under which latent representations acquire symbolic stability (Volume II), and how these structures govern large-scale computational organization (Volume III).
